\scanpage{30}
$Z$ est un diviseur sur $X\times T$ et $\unsure{\mathfrak{a}|\alpha|a|\mathfrak{t}}$ sur $T$ est de degré $0$ et tel que $S_T(\mathfrak{a})=0$, alors $Z(\mathfrak{a})$ est lin. équiv. à $0$. Or $\unsure{c^1|c'|\text{on a}}Z(\mathfrak{a})=u(S_T(\mathfrak{a}))$, où $u:\deleted{\mathrm{Alb}}\;A_T\longrightarrow\underline{\mathrm{Pic}}_X^0$ est l'hom défini par $Z$.

\uline{Théorème} 1.\deleted{\unsure{5}}6. Soient $z,z'$ deux classes dans $A'{}^*(X)$ alg. équiv. à $0$, $X$ étant \deleted{[projective]}\inserted{qu. pr.}. Alors $zz'=0$. [Si $X$ est projective, alors il suffit que l'un des \deleted{facteurs} $z$ soit alg. équiv. à $0$, et l'autre \uline{num.} équiv. à $0$ ??]

% slanted
\margin{à voir \unsure{si il|s'il} ne \unsure{suffirait|suffit} pas que $z$ ou $z'$ ne soit que \unsure{$\tau$-équiv.|$z$-équiv.|num. équiv.} à $0$, voire (si $X$ \unsure{projectif|projec.}) équiv. \uline{num.} à $0$.}

\deleted{Cela}\inserted{La première assertion} résulte formellement de la définition de ``alg. équiv. à $0$'' \inserted{via variétés $T$ \uline{projectives} \uline{lisses} connexes,} et du fait que si $T,T'$ sont deux \inserted{telles} variétés \deleted{et} (lisse connexe suffit) et $x,y$ sont deux points de $T$, $x'$ et $y'$ deux pts de $T'$, alors $\mathfrak{a}=\inserted{((x)-(y))\times((x')-(y'))=}(x,x')-(x,y')-(y,x')+(y,y')$ est tel que $\deleted{\unsure{S_T(\mathfrak{a})|S_{\mathfrak{a}}}}\;S_{T\times T'}(\mathfrak{a})=0$.

\margin{des courbes si on veut !}

% slanted
\margin{NB. On ignore si cet $\mathfrak{a}$ est lin. équiv. à $0$, $\overline{\mathrm{m}}$ si $T,T'$ sont des \uline{courbes} !}

[Il y a intérêt à transformer le \deleted{Pb}\inserted{l'énoncé} en un énoncé sur l'opération $\times$]

% archive sheet number: 29
